%
% PropagandaTrace — CS 544 NLP Midterm Report
% ACL 2023 Format (acl.sty required — download from ACL Anthology template)
% Compile: pdflatex midterm_report.tex  (twice for references)
%
\documentclass[11pt]{article}
\usepackage[hyperref]{acl}      % acl.sty from ACL 2023/2024 template package
\usepackage{times}
\usepackage{latexsym}
\usepackage[T1]{fontenc}
\usepackage[utf8]{inputenc}
\usepackage{microtype}
\usepackage{inconsolata}
\usepackage{booktabs}
\usepackage{amsmath}
\usepackage{graphicx}
\usepackage{url}
\usepackage{multirow}

\title{PropagandaTrace: Midterm Progress Report\\
Multi-Source Attribution of AI-Generated Wartime Propaganda\\via Watermark Fingerprinting}

\author{
  Samarth Rajendra$^{1}$\quad
  Sanath Nagendra Bhargav$^{1}$\quad
  Jeevika Kiran$^{1}$ \\
  \textbf{Shree Sriadibhatla}$^{1}$\quad
  \textbf{Akash Gandi}$^{1}$ \\[4pt]
  $^{1}$University of Southern California \\
  \texttt{\{samarth.rajendra, sanathna, jeevikak, sriadibh, gandi\}@usc.edu}
}

\begin{document}
\maketitle

% ─────────────────────────────────────────────────────────────────────────────
\section{Report on Tasks Performed}
% ─────────────────────────────────────────────────────────────────────────────

\subsection{Project Overview}
PropagandaTrace addresses a critical gap in AI safety: existing methods detect
\emph{whether} text is AI-generated, but cannot identify \emph{which} model
produced it~\cite{kirchenbauer2023watermark,sadasivan2023aigenerated}.
Our framework assigns each candidate language model a unique watermark key at
deployment time and performs forensic multi-class attribution by comparing
recovered text against each key's signal.

\subsection{Data Collection (Phase 1)}

We built a four-source seed-prompt corpus for propagating through the generation
pipeline.  First, we downloaded the \textbf{WTWT} dataset~\cite{conforti2020wtwt}
via the HuggingFace \texttt{datasets} library, which provides 4{,}000 stance-labeled
tweets originally covering merger/acquisition rumours---repurposed here as
topically neutral trigger texts whose surface framing is overridden by our
propaganda prompt templates.  Second, we queried the \textbf{GDELT} v2 DocAPI
(no authentication required) with six conflict-oriented keyword strings
(\textit{``military offensive attack,''} \textit{``war propaganda
disinformation,''} etc.) to retrieve up to 500 recent headline–snippet pairs
covering documented conflict zones.  Third, we fetched article paragraphs from
the public development split of \textbf{SemEval-2020 Task~11}~\cite{semeval2020t11},
which provides professionally annotated propaganda fragments.  Finally, we
implemented a template-based prompt generator containing 20 fill-in-the-blank
templates (10 military situation reports; 10 inflammatory political
commentaries) populated from curated word lists covering factions, adversaries,
and geographic references drawn from documented disinformation patterns.  After
deduplication, the pipeline produced \textbf{4{,}000 seed prompts} stored in a
JSONL file with source provenance metadata.

\subsection{Watermarking Implementation (Phase 2)}

We implemented both target watermarking schemes as HuggingFace
\texttt{LogitsProcessor} objects, which plug directly into
\texttt{model.generate()} without any model fine-tuning.

\paragraph{KGW Scheme.}
Following~\citet{kirchenbauer2023watermark}, at each generation step we hash
the previous token together with a model-specific secret key
($\texttt{SHA-256}(\text{key} \| t_{i-1})$) to seed a pseudorandom permutation
of the vocabulary.  The first $\gamma|V|$ tokens in the permutation form the
\emph{green list}; their logits receive an additive bias $\delta$.  We set
$\gamma = 0.25$ and $\delta = 2.0$ following the paper's recommended defaults.
Detection computes the z-score
$z = (|\text{green}| - \gamma T) / \sqrt{\gamma(1-\gamma)T}$
over a tokenised text of length $T$; we flag a text as watermarked at
$z > 4.0$.  Each of the three models (Mistral-7B, LLaMA-3-8B, Falcon-7B) uses
a distinct key, enabling multi-class attribution.

\paragraph{Aaronson / EXP Scheme.}
We implement the exponential minimum sampling scheme~\cite{aaronson2022watermark}
via the Gumbel-max trick: at position $t$ we derive a uniform random vector
$\mathbf{r}_t \in [0,1]^{|V|}$ from $(\text{key}, t)$ and replace logits with
$\log p_i - \log(-\log r_{t,i})$.  Taking the argmax is equivalent to sampling
$\text{argmax}_i\, r_{t,i}^{1/p_i}$, which is skewed toward tokens $i$ where
$r_{t,i}$ is large.  Detection accumulates
$s_t = -\log(1 - r_{t,w_t})$ over observed tokens $w_t$.

\paragraph{Corpus Generation Status.}
We configured all three LLMs to load in 4-bit NF4 quantization via
\texttt{bitsandbytes}~\cite{dettmers2022llmint8}, reducing per-model VRAM to
approximately 6\,GB and enabling sequential execution on a single A100.  As of
the midterm submission, we have generated \textbf{3{,}000 texts from
Mistral-7B-Instruct-v0.2} under the KGW scheme (one-third of the primary KGW
corpus), and have verified end-to-end detection with $z$-scores averaging $6.8$
on unattacked watermarked text.  Falcon-7B and LLaMA-3-8B generation is
scheduled to complete within the next two weeks using USC's CARC HPC cluster.

\subsection{Evasion Pipeline (Phase 3 — Framework Complete)}

We have implemented all four planned evasion modules as standalone Python
scripts:

\begin{itemize}
  \item \textbf{T5 Lexical Paraphrase}: single-pass rewriting via T5-large with
    a \texttt{``paraphrase:''} prefix.
  \item \textbf{PEGASUS Neural Paraphrase}: aggressive structural rewriting via
    PEGASUS-XSUM~\cite{zhang2020pegasus}, which was originally trained for
    extreme summarisation and produces fluent yet structurally distant rewrites.
  \item \textbf{Round-Trip Translation (Arabic)}: English $\to$ Arabic $\to$
    English via NLLB-200~\cite{nllbteam2022nllb}.
  \item \textbf{Round-Trip Translation (Russian)}: English $\to$ Russian $\to$
    English via NLLB-200.
\end{itemize}

Each evasion module reads a Phase~2 JSONL corpus file, appends the evaded text
alongside the original, and writes a new JSONL, making the output suitable for
direct watermark detection and attribution scoring in Phase~4.

% ─────────────────────────────────────────────────────────────────────────────
\section{Potential Risks and Challenges}
% ─────────────────────────────────────────────────────────────────────────────

\paragraph{R1 — Compute Availability.}
Generating 9{,}000 texts across three 7–8\,B parameter models is time-intensive
even with 4-bit quantization.  Each model requires $\sim$6\,GB VRAM and
produces roughly 30--50 texts per minute on an A100.  Queue delays on shared HPC
nodes could compress our implementation schedule.

\paragraph{R2 — Watermark Fragility Under Aggressive Evasion.}
PEGASUS and round-trip translation substantially alter token sequences; there is
a real risk that the watermark z-score falls below the detection threshold for a
majority of attacked texts, making attribution impossible without a separate
semantic fingerprint.  Preliminary literature suggests 40--60\% survival rates
for KGW under paraphrase~\cite{kirchenbauer2023watermark}.

\paragraph{R3 — Attribution Ambiguity Between Similar Models.}
Mistral-7B and LLaMA-3-8B share substantial vocabulary overlap and similar
pretraining distributions.  Their green-list distributions under different keys
may still be correlated, inflating the false attribution rate when the
confidence margin between the top two scores is small.

\paragraph{R4 — Gated Model Access.}
Meta-Llama-3-8B requires accepting a usage licence on the HuggingFace Hub.
Licence approval for all team members is pending, which could delay LLaMA corpus
generation by 1--2 weeks.

\paragraph{R5 — Dataset Quality and Propaganda Realism.}
Template-generated prompts may produce formulaic text that does not reflect
genuine wartime disinformation style, reducing ecological validity and inflating
detection performance relative to real-world deployment.

\paragraph{R6 — Evaluation Metric Calibration.}
The z-score threshold for KGW detection ($z > 4.0$) and the Aaronson mean-score
threshold ($> 0.8$) were chosen from prior work on non-propaganda text; these
may need recalibration for our specific generation distribution and text length
range (150--300 tokens).

% ─────────────────────────────────────────────────────────────────────────────
\section{Risk Mitigation Plan}
% ─────────────────────────────────────────────────────────────────────────────

\paragraph{Compute (R1).}
We will submit parallel SLURM jobs on CARC, one per model, using our
\texttt{phase2\_generate\_corpus.py --model} flag.  Our generation script
supports mid-run resumption via line-count checkpointing, so interrupted jobs
do not require restarting.  As a fallback, we will use the HuggingFace
Inference API for Falcon-7B, which does not require local GPU.

\paragraph{Watermark Fragility (R2).}
We will tune $\delta$ in $\{1.5, 2.0, 3.0\}$ and $\gamma$ in $\{0.25, 0.50\}$
to find configurations that survive evasion while maintaining text quality.  We
will also implement a \emph{hybrid scoring} mode that combines KGW z-score with
Aaronson mean-score, and fall back to a TF-IDF stylometric baseline if both
watermarks are destroyed.

\paragraph{Attribution Ambiguity (R3).}
The fallback linear classifier over the three watermark scores (one per key)
will be trained with a reject-option threshold: if no key's score clears a
confidence margin of 1.5 z-points above the second-best, the system abstains
rather than misattributing.

\paragraph{Gated Model Access (R4).}
We have already initiated licence requests for LLaMA-3-8B on all five team
accounts.  If approval is delayed, we will substitute
\texttt{meta-llama/Llama-2-7b-chat-hf}, which shares the same tokeniser and
architecture family.

\paragraph{Dataset Quality (R5).}
We will add 200 real propaganda excerpts from the SemEval-2020 Task~11 articles
as seed material and manually review 50 generated texts per model to verify
stylistic plausibility.

\paragraph{Threshold Calibration (R6).}
We will hold out 10\% of the KGW corpus as a calibration set and sweep
thresholds to achieve $\le 5\%$ false positive rate before reporting final
detection metrics.

% ─────────────────────────────────────────────────────────────────────────────
\section{Individual Contributions}
% ─────────────────────────────────────────────────────────────────────────────

\begin{table}[h]
\centering
\small
\begin{tabular}{@{}lp{5.5cm}@{}}
\toprule
\textbf{Member} & \textbf{Contribution} \\
\midrule
Samarth Rajendra
  & Implemented the KGW \texttt{LogitsProcessor} and z-score detector;
    ran end-to-end generation experiments on Mistral-7B; verified
    $z$-score calibration. \\[4pt]
Sanath Nagendra Bhargav
  & Implemented the Aaronson (EXP) \texttt{LogitsProcessor} and
    mean-score detector; authored \texttt{config.yaml} and
    generation pipeline architecture. \\[4pt]
Jeevika Kiran
  & Designed and implemented the data collection pipeline (WTWT,
    GDELT, SemEval-2020 T11); created propaganda prompt templates;
    ran Phase~1 and curated seed prompts. \\[4pt]
Shree Sriadibhatla
  & Built \texttt{phase2\_generate\_corpus.py} with 4-bit quantization
    via \texttt{bitsandbytes}; implemented SLURM job scripts for CARC;
    debugged tokeniser padding for all three models. \\[4pt]
Akash Gandi
  & Implemented all four evasion modules (T5, PEGASUS,
    NLLB-200 Arabic, NLLB-200 Russian); authored
    \texttt{phase3\_evasion.py}; established
    output schema for Phase~4 attribution. \\
\bottomrule
\end{tabular}
\caption{Individual team contributions at the midterm milestone.}
\label{tab:contributions}
\end{table}

% ─────────────────────────────────────────────────────────────────────────────
\bibliography{references}
% ─────────────────────────────────────────────────────────────────────────────

\end{document}
